% ================================================================
\chapter{Heaps binomiais, heaps de Fibonacci e análise amortizada}\label{cap:heapav}\label{sec:heapav}
% ================================================================
\begin{objetivos}
\begin{itemize}
  \item Árvores binomiais e o heap binomial, cuja fusão imita a soma de números binários.
  \item A análise amortizada pelos métodos do banqueiro e do potencial.
  \item O heap de Fibonacci: a ideia de adiar trabalho, as marcas e os cortes em cascata, e por que o grau dos nós é logarítmico.
  \item Como escolher entre os heaps binário, binomial e de Fibonacci.
\end{itemize}
\end{objetivos}

O heap binário do capítulo anterior tem um ponto fraco: \textbf{unir} duas listas de prioridades exige concatenar os vetores e reconstruir o heap, o que custa $\Theta(n)$. Imagine um sistema operacional com uma fila de processos por processador: quando um processador é desligado, a sua fila precisa ser fundida com a de outro. Em aplicações com muitas fusões, pagar $\Theta(n)$ a cada uma é inviável.

Este capítulo apresenta duas estruturas em que a fusão é barata. O \textbf{heap binomial} representa a lista de prioridades como uma coleção de árvores cujas ordens correspondem aos bits de $n$, e funde duas listas como quem soma dois números binários, em $O(\log n)$ no pior caso. O \textbf{heap de Fibonacci} vai além: adia toda a arrumação para o momento da remoção do máximo, e só é eficiente quando se considera o custo de uma \emph{sequência} de operações. Por isso a análise amortizada, que formaliza esse tipo de garantia, aparece aqui.

% ------------------------------------------------------------------
\section{Árvores binomiais}\label{sec:binomial}
\begin{defi}
A árvore binomial $B_0$ é um único nó. Para $k\ge1$, $B_k$ é obtida \textbf{subordinando} (operação \emph{link}) uma $B_{k-1}$ a outra $B_{k-1}$: a raiz de uma passa a ser filha da raiz da outra. Num heap, a raiz de \emph{maior} prioridade fica por cima.
\end{defi}

\begin{figura}[As árvores binomiais $B_0$ a $B_4$. Em cada $B_k$, a subárvore em destaque é a cópia de $B_{k-1}$ que foi subordinada à outra.\label{fig:binomiais}]
\forestset{bn/.style={for tree={circle, draw, minimum size=2.8mm, inner sep=0pt, s sep=1.6mm, l sep=1.6mm, edge={thick}}}}
\begin{tabular}{@{}c@{\hspace{7mm}}c@{\hspace{7mm}}c@{\hspace{7mm}}c@{\hspace{7mm}}c@{}}
\begin{forest} bn, baseline [] \end{forest} &
\begin{forest} bn, baseline [ [,fill=destaque] ] \end{forest} &
\begin{forest} bn, baseline [ [,fill=destaque [,fill=destaque] ] [] ] \end{forest} &
\begin{forest} bn, baseline [ [,fill=destaque [,fill=destaque [,fill=destaque]] [,fill=destaque]] [ [] ] [] ] \end{forest} &
\begin{forest} bn, baseline [ [,fill=destaque [,fill=destaque [,fill=destaque [,fill=destaque]] [,fill=destaque]] [,fill=destaque [,fill=destaque]] [,fill=destaque]] [ [ [] ] [] ] [ [] ] [] ] \end{forest}\\[2pt]
$B_0$ & $B_1$ & $B_2$ & $B_3$ & $B_4$
\end{tabular}
\end{figura}

\begin{prop}\label{prop:binom}
Para todo $k\ge0$:
\begin{enumerate}
  \item $B_k$ tem exatamente $2^k$ nós;
  \item o grau da raiz, chamado \textbf{ordem} de $B_k$, é $k$, e é o maior grau da árvore;
  \item a altura de $B_k$, medida em arestas, é $k$;
  \item há exatamente $\binom kj$ nós na profundidade $j$;
  \item os filhos da raiz de $B_k$ são raízes de árvores $B_{k-1},B_{k-2},\dots,B_1,B_0$.
\end{enumerate}
\end{prop}
\begin{proof}
Todas as propriedades seguem por indução em $k$, olhando para o \emph{link} que forma $B_k$ a partir de duas cópias de $B_{k-1}$: o número de nós dobra ($2\cdot2^{k-1}$); a raiz ganha um filho; a árvore subordinada fica um nível abaixo, e a altura passa a $\max\{k-1,\,1+(k-1)\}=k$; os nós na profundidade $j$ vêm da cópia de cima (profundidade $j$) e da de baixo (profundidade $j-1$), e $\binom{k-1}{j}+\binom{k-1}{j-1}=\binom kj$ --- daí o nome ``binomial''; e o novo filho da raiz é uma $B_{k-1}$, que se junta aos filhos $B_{k-2},\dots,B_0$ que a raiz já tinha.
\end{proof}

\begin{exemplo}[Conferindo as propriedades em $B_3$]
Olhe para $B_3$ na Figura~\ref{fig:binomiais}. Ele tem $8=2^3$ nós; a raiz tem 3 filhos e nenhum outro nó tem grau maior; o caminho mais longo da raiz até uma folha tem 3 arestas. Contando os nós por profundidade, obtemos $1,3,3,1$, que são os coeficientes $\binom30,\binom31,\binom32,\binom33$ --- a linha 3 do triângulo de Pascal. Por fim, os filhos da raiz, da esquerda para a direita, são raízes de uma $B_2$ (4 nós), uma $B_1$ (2 nós) e uma $B_0$ (1 nó).
\end{exemplo}

\begin{cuidado}[Duas convenções de altura]
A altura de $B_k$ é $k$ quando medida em arestas, e $k+1$ quando medida em níveis (com folhas de altura 1, como nos heaps binários do Capítulo~\ref{cap:heap}). É a mesma árvore. Em exercícios sobre alturas, declare a convenção usada. Neste livro, a ordem de $B_k$ é o grau da raiz, $k$, e a altura de $B_k$ é $k$ arestas.
\end{cuidado}

% ------------------------------------------------------------------
\section{Heap binomial}\label{sec:heapbinomial}
\begin{defi}
Um \textbf{heap binomial} é uma floresta $H$ de árvores binomiais em que (1) cada árvore satisfaz a condição de heap (filho $\le$ pai) e (2) as árvores têm \textbf{ordens distintas}, isto é, há no máximo uma árvore de cada ordem. As raízes ficam numa lista, em ordem crescente de ordem.
\end{defi}

\begin{prop}
(a) Para cada $n$ existe um único formato de heap binomial com $n$ nós: $B_k$ pertence a $H$ se e somente se o bit $k$ da representação binária de $n$ é 1. (b) Portanto $H$ tem no máximo $\floor{\log_2n}+1$ árvores, cada uma de altura $\le\floor{\log_2n}$.
\end{prop}
\begin{proof}
A árvore $B_k$ contribui com $2^k$ nós e as ordens são distintas, logo $n=\sum_{B_k\in H}2^k$: é a representação binária de $n$, que é única. Um número $n$ tem $\floor{\log_2n}+1$ bits.
\end{proof}

\begin{exemplo}[Um heap binomial com 13 elementos]
Como $13=1101_2=8+4+1$, um heap binomial com 13 elementos tem exatamente as árvores $B_0$, $B_2$ e $B_3$ (Figura~\ref{fig:binheap13}). Do mesmo modo, $26=11010_2$ dá $H=\{B_1,B_3,B_4\}$. O máximo do heap é a maior das raízes, aqui o 50.
\begin{figura}[Um heap binomial com $13=1101_2$ elementos. Cada árvore é um heap (pai $\ge$ filhos), e há uma árvore para cada bit 1 de $n$. A lista de raízes está em ordem crescente de ordem.\label{fig:binheap13}]
\begin{tabular}{@{}c@{\hspace{10mm}}c@{\hspace{10mm}}c@{}}
\begin{forest} bstp, baseline [12] \end{forest} &
\begin{forest} bstp, baseline [41 [33 [17]] [28]] \end{forest} &
\begin{forest} bstp, baseline [50,feito [44 [31 [10]] [25]] [37 [29]] [6]] \end{forest}\\[2pt]
\small $B_0$ (bit 0) & \small $B_2$ (bit 2) & \small $B_3$ (bit 3)
\end{tabular}\\[4pt]
{\small\color{MidnightBlue} $13=\underbrace{1}_{B_3}\underbrace{1}_{B_2}\underbrace{0}_{\text{sem }B_1}\underbrace{1}_{B_0}{}_2$}
\end{figura}
\end{exemplo}

\subsection{Operações}
A operação central é a \textbf{fusão}, e as demais são escritas em termos dela.
\begin{pseudo}{Link e fusão de heaps binomiais}
\begin{algorithmic}
\Procedure{Link}{$y,z$} \Comment{$y$ e $z$ são raízes de árvores $B_{k-1}$; $z$ fica por cima}
  \State $pai[y]\gets z$;\quad $\textit{irmão}[y]\gets filho[z]$;\quad $filho[z]\gets y$;\quad $grau[z]\gets grau[z]+1$
\EndProcedure
\Procedure{Fusão}{$H',H''$}
  \State $H := H'\cup H''$ \Comment{intercala as listas de raízes por ordem crescente: $O(\log n)$}
  \While{$H$ contém duas árvores $B,B'$ de mesma ordem} \Comment{menor ordem primeiro}
    \If{$B.raiz.chave \ge B'.raiz.chave$} \State subordinar $B'$ a $B$
    \Else \State subordinar $B$ a $B'$
    \EndIf
  \EndWhile
\EndProcedure
\end{algorithmic}
\end{pseudo}

\begin{intuicao}[Fundir é somar em binário]
Na soma $011_2+011_2$, a coluna das unidades tem dois 1, que viram um 0 e um ``vai um'' para a coluna seguinte. Na fusão, duas árvores $B_0$ viram uma $B_1$, que ``vai'' para a posição das árvores de ordem 1. Ali podem se encontrar até três árvores de mesma ordem (duas originais e a que veio da posição anterior): uma fica, e as outras duas se juntam e vão para a posição seguinte, exatamente como $1+1+1=11_2$.
\end{intuicao}

Há $O(\log n)$ posições e, em cada uma, no máximo três árvores de mesma ordem ao mesmo tempo, de modo que a fusão custa $O(\log n)$ no pior caso. Numa implementação com ponteiros, a lista de raízes é percorrida mantendo o nó anterior, o atual e o próximo, e há quatro casos: ordens diferentes, três árvores seguidas de mesma ordem, e duas árvores de mesma ordem em que vence a atual ou a seguinte.

\begin{exemplo}[Fusão passo a passo]\label{ex:fusaobin}
Sejam $H_1=\{7,\ 9\to3\}$ e $H_2=\{5,\ 8\to6\}$, cada um com $3=11_2$ elementos: uma $B_0$ e uma $B_1$ (Figura~\ref{fig:fusaobin}).
\begin{enumerate}
  \item Ordem 0: há duas $B_0$, com raízes 7 e 5. O link põe a maior por cima e forma a $B_1$ $7\to5$ (o ``vai um'').
  \item Ordem 1: agora há três $B_1$: $7\to5$, $9\to3$ e $8\to6$. A primeira fica; as outras duas se juntam, com 9 por cima, numa $B_2$.
  \item O resultado é $\{B_1, B_2\}$, com $6=110_2$ elementos: de fato, $11_2+11_2=110_2$.
\end{enumerate}
\begin{figura}[A fusão de $H_1$ com $H_2$ (Exemplo~\ref{ex:fusaobin}) imita a soma $11_2+11_2=110_2$.\label{fig:fusaobin}]
\begin{tabular}{@{}c@{\hspace{9mm}}c@{\hspace{9mm}}c@{}}
\begin{forest} bstp, baseline [7] \end{forest}\ \begin{forest} bstp, baseline [9 [3]] \end{forest}
\ \ \raisebox{0.5ex}{\small$+$}\ \
\begin{forest} bstp, baseline [5] \end{forest}\ \begin{forest} bstp, baseline [8 [6]] \end{forest}
&
\begin{forest} bstp, baseline [7,novo [5]] \end{forest}\ \begin{forest} bstp, baseline [9 [3]] \end{forest}\ \begin{forest} bstp, baseline [8 [6]] \end{forest}
&
\begin{forest} bstp, baseline [7 [5]] \end{forest}\ \begin{forest} bstp, baseline [9,novo [8 [6]] [3]] \end{forest}\\[3pt]
\small (a) $H_1$ e $H_2$ & \small (b) $B_0+B_0\to B_1$ & \small (c) duas das três $B_1$ viram $B_2$
\end{tabular}
\end{figura}
\end{exemplo}

\begin{exemplo}[{Fusão de $H_1=\{B_1,B_2,B_3\}$ com $H_2=\{B_0,B_1,B_3\}$}]\label{ex:fusaoformas}
Olhando apenas para os formatos: os heaps têm $14=1110_2$ e $11=1011_2$ elementos. A $B_0$ de $H_2$ fica sozinha na ordem 0. Na ordem 1, as duas $B_1$ formam uma $B_2$ (o ``vai um''); na ordem 2, essa $B_2$ se junta à $B_2$ de $H_1$ e forma uma $B_3$; na ordem 3 há, então, três árvores --- as $B_3$ de $H_1$ e de $H_2$ e a que acabou de se formar ---, e, como na soma $1+1+1=11_2$, uma delas fica e as outras duas formam uma $B_4$. São três links, e o resultado é $\{B_0,B_3,B_4\}$, isto é, $25=11001_2=1110_2+1011_2$.
\end{exemplo}

As demais operações são as seguintes.
\begin{itemize}
  \item \textbf{Inserção}: fusão com um heap $\{B_0\}$ que contém só o novo elemento. No pior caso custa $O(\log n)$, quando $n=2^k-1$, todos os bits são 1 e há $k$ links. Mas inserir equivale a somar 1 a um contador binário, e cada link corresponde a um ``vai um''. Se os bits de $n$ são equiprováveis, a inserção faz em média um link, pois a posição $j$ do primeiro 0 de $n$ (lido da direita para a esquerda) tem esperança $\sum_{j\ge1}j/2^j=2$. E, ao longo de uma sequência de inserções, o custo \textbf{amortizado} é $O(1)$, como mostra o exemplo do contador binário da §\ref{sec:amort}.
  \item \textbf{Seleção do máximo}: o máximo é uma das no máximo $\floor{\log n}+1$ raízes, e é encontrado em $O(\log n)$, ou em $O(1)$ se um ponteiro para a maior raiz for mantido.
  \item \textbf{Remoção do máximo}: encontre a raiz máxima $r$, numa árvore $B_k$, e retire $B_k$ de $H$, formando $H'=H-B_k$. Removendo $r$, seus filhos $B_{k-1},\dots,B_0$ formam um heap binomial $H''$. Faça a fusão $H'\cup H''$. Total: $O(\log n)$.
  \item \textbf{Aumento de prioridade}: o nó sobe dentro da própria árvore trocando de lugar com o pai, como no heap binário, em $O(\log n)$, pois a altura é $\le\log n$.
  \item \textbf{Diminuição de prioridade}: remova o nó (veja abaixo) e reinsira-o com a nova prioridade, em $O(\log n)$. Fazê-lo descer trocando com o maior filho também funciona, mas cada nível tem até $O(\log n)$ filhos a comparar, o que dá $O(\log^2n)$.
  \item \textbf{Remoção de um nó qualquer}: aumente sua prioridade para $+\infty$, para que ele suba até a raiz, e remova o máximo: $O(\log n)$.
  \item \textbf{Construção}: $n$ inserções sucessivas custam $O(n)$ no total, pelo mesmo argumento do contador binário.
\end{itemize}

\begin{exemplo}[Remoção do máximo]\label{ex:rembin}
No heap resultante do Exemplo~\ref{ex:fusaobin}, $\{7\to5,\ B_2 \text{ de raiz } 9\}$, o máximo é a raiz 9. Retirando a sua árvore, sobra $H'=\{7\to5\}$. Os filhos de 9 formam outro heap binomial, $H''=\{3,\ 8\to6\}$ (uma $B_0$ e uma $B_1$). A fusão $H'\cup H''$ encontra duas $B_1$, $7\to5$ e $8\to6$, e as junta com 8 por cima. O resultado tem $5=101_2$ elementos (Figura~\ref{fig:rembin}).
\begin{figura}[Remoção do máximo: os filhos da raiz removida formam um heap binomial, que é fundido com o resto.\label{fig:rembin}]
\begin{tabular}{@{}c@{\hspace{10mm}}c@{\hspace{10mm}}c@{}}
\begin{forest} bstp, baseline [7 [5]] \end{forest}\ \begin{forest} bstp, baseline [9,desreg [8 [6]] [3]] \end{forest}
&
\begin{forest} bstp, baseline [7 [5]] \end{forest}\ \ \raisebox{0.5ex}{\small$\cup$}\ \ \begin{forest} bstp, baseline [3] \end{forest}\ \begin{forest} bstp, baseline [8 [6]] \end{forest}
&
\begin{forest} bstp, baseline [3] \end{forest}\ \begin{forest} bstp, baseline [8,novo [7 [5]] [6]] \end{forest}\\[3pt]
\small (a) o máximo é 9 & \small (b) $H'$ e os filhos de 9 & \small (c) após a fusão
\end{tabular}
\end{figura}
\end{exemplo}

% ------------------------------------------------------------------
\section{Análise amortizada}\label{sec:amort}
\begin{intuicao}
Pense num cofrinho. Toda vez que você faz algo barato, paga um pouco a mais e guarda o troco no cofrinho. Quando chega uma despesa grande, ela é paga com o dinheiro guardado. Se o cofrinho nunca fica negativo, o total gasto nunca passou do que você ``cobrou'' de cada operação. A análise amortizada é a contabilidade desse cofrinho.
\end{intuicao}

Em várias estruturas, uma operação individual pode ser cara, mas operações caras são raras, porque cada uma delas desfaz um trabalho acumulado por muitas operações baratas. A \textbf{análise amortizada} mede o custo total de uma \emph{sequência} de operações, no pior caso, e o divide pelo número de operações. É uma garantia de pior caso para a sequência e \textbf{não envolve probabilidade}.

\paragraph{Método do banqueiro (créditos).} Cada operação recebe um número fixo de créditos, que é o seu custo amortizado. Operações baratas deixam créditos guardados na estrutura, e esses créditos pagam as operações caras mais tarde. Se o saldo de créditos nunca fica negativo, o custo real total é no máximo a soma dos custos amortizados.

\paragraph{Método do potencial (físico).} Escolhe-se uma função $\Phi\ge0$ sobre os estados da estrutura, que mede o trabalho ``armazenado'' (o dinheiro no cofrinho). O custo amortizado da $i$-ésima operação, de custo real $t_i$, é $a_i=t_i+\Phi(E_i)-\Phi(E_{i-1})$. Somando, os termos de potencial se cancelam:
\[
\sum_{i=1}^m t_i=\sum_{i=1}^m a_i-\Phi_m+\Phi_0\ \le\ \sum_{i=1}^m a_i\qquad\text{sempre que }\Phi_m\ge\Phi_0 .
\]

\begin{exemplo}[Contador binário]\label{ex:contador}
Incrementar um contador de $k$ bits pode trocar $k$ bits (de $0111$ para $1000$, por exemplo), mas isso é raro. A tabela mostra os oito primeiros incrementos a partir de zero:
\begin{center}\small
\begin{tabular}{@{}rcccccc@{}}
\toprule
incremento & contador & bits trocados ($t_i$) & $\Phi$ (bits 1) & $\Delta\Phi$ & $a_i=t_i+\Delta\Phi$ & total real\\
\midrule
--- & \texttt{0000} & --- & 0 & --- & --- & 0\\
1 & \texttt{000\textbf{1}} & 1 & 1 & $+1$ & 2 & 1\\
2 & \texttt{00\textbf{10}} & 2 & 1 & $0$ & 2 & 3\\
3 & \texttt{001\textbf{1}} & 1 & 2 & $+1$ & 2 & 4\\
4 & \texttt{0\textbf{100}} & 3 & 1 & $-1$ & 2 & 7\\
5 & \texttt{010\textbf{1}} & 1 & 2 & $+1$ & 2 & 8\\
6 & \texttt{01\textbf{10}} & 2 & 2 & $0$ & 2 & 10\\
7 & \texttt{011\textbf{1}} & 1 & 3 & $+1$ & 2 & 11\\
8 & \texttt{\textbf{1000}} & 4 & 1 & $-2$ & 2 & 15\\
\bottomrule
\end{tabular}
\end{center}
Os custos reais oscilam (1, 2, 1, 3, 1, 2, 1, 4), mas o custo amortizado é sempre 2. Em geral, com $\Phi=$ número de bits iguais a 1, um incremento que zera $t$ bits e liga um tem custo amortizado $a=(t+1)+(1-t)=2$. Logo $m$ incrementos trocam no máximo $2m$ bits: na tabela, $15\le16$. É o mesmo argumento que mostra que as inserções num heap binomial custam $O(1)$ amortizado: inserir é somar 1.
\end{exemplo}

\begin{exemplo}[Pilha com desempilhamento múltiplo]\label{ex:pilha}
Cada operação desempilha $k\ge0$ elementos e depois empilha um; uma operação isolada pode custar $\Theta(n)$. Com $\Phi=$ número de elementos na pilha, uma operação que encontra $p$ elementos tem custo amortizado $a=(1+k)+(p-k+1)-p=2$. Logo $m$ operações custam no máximo $2m=O(m)$.

Pelo método do banqueiro, a mesma conta fica assim: cada operação paga 2 créditos; um paga o empilhamento e o outro fica ``grudado'' no elemento empilhado, para pagar o seu futuro desempilhamento. Na sequência com $k=0,0,0,3,0$, os custos reais são $1,1,1,4,1$, o total é $8$, e os créditos pagos somam $2\cdot5=10$. Na quarta operação, os três créditos guardados nos elementos da pilha pagam os três desempilhamentos.
\end{exemplo}

\begin{cuidado}[Amortizado não é caso médio]
Uma garantia amortizada vale para \emph{toda} sequência de operações, sem hipóteses probabilísticas, mas não limita o custo de uma operação \emph{individual}, que ainda pode ser alto (o oitavo incremento da tabela custou 4). Isso é decisivo quando um sistema exige que cada operação termine num tempo previsível.
\end{cuidado}

% ------------------------------------------------------------------
\section{Heap de Fibonacci}\label{sec:fib}
O heap binomial é ``rigoroso'': cada inserção ou fusão resolve imediatamente as colisões de ordem. O heap de Fibonacci é \textbf{preguiçoso}: adia toda a arrumação até a remoção do máximo, quando ela é feita de uma só vez.

\begin{intuicao}
É como deixar a louça acumular na pia e lavar tudo de uma vez. Cada refeição (inserção) é rápida; a lavagem (consolidação) é demorada, mas só acontece de vez em quando e é paga pelas refeições que sujaram a louça. A remoção do máximo é o momento da lavagem: como é preciso achar o novo máximo de qualquer jeito, aproveita-se para arrumar a estrutura.
\end{intuicao}

\paragraph{Estrutura.} É uma floresta de árvores com a condição de heap. As raízes formam uma \textbf{lista circular duplamente encadeada}, e os filhos de cada nó também. Um ponteiro indica a raiz máxima. Cada nó guarda a chave, um ponteiro para o pai, um para um filho, dois para os irmãos esquerdo e direito, o grau e um bit de \textbf{marca}. As listas circulares duplas permitem inserir ou retirar um nó e concatenar duas listas em $O(1)$.

\begin{figura}[Um heap de Fibonacci. As raízes formam uma lista circular dupla (setas duplas; a última volta para a primeira) e o ponteiro \textsf{max} aponta a maior raiz. Os nós escuros estão marcados (já perderam um filho). Os filhos de cada nó também formam uma lista circular dupla, omitida no desenho.\label{fig:fibestrutura}]
\begin{tikzpicture}[x=1.25cm,y=0.95cm, every node/.append style={font=\scriptsize}]
\node[vtxs] (r1) at (0,0) {23};
\node[vtxs] (r2) at (1.3,0) {7};
\node[vtxs, fill=feito] (r3) at (2.6,0) {48};
\node[vtxs] (r4) at (4.4,0) {17};
\node[vtxs] (r5) at (5.7,0) {24};
\foreach \a/\b in {r1/r2,r2/r3,r3/r4,r4/r5} {\draw[{Stealth[length=1.5mm]}-{Stealth[length=1.5mm]}, gray] (\a) -- (\b);}
\draw[-{Stealth[length=1.5mm]}, gray, rounded corners] (r5.north) -- ++(0,0.45) -| (r1.north);
\node[font=\scriptsize\sffamily, MidnightBlue] (mx) at (2.6,1.25) {max};
\draw[seta, MidnightBlue] (mx) -- (r3);
\node[vtxs, fill=black!70, text=white] (a1) at (2.1,-1) {39};
\node[vtxs] (a2) at (3.1,-1) {41};
\node[vtxs, fill=black!70, text=white] (a3) at (2.1,-2) {18};
\node[vtxs] (a4) at (2.1,-3) {3};
\node[vtxs] (a5) at (3.1,-2) {38};
\draw[aresta] (r3)--(a1) (r3)--(a2) (a1)--(a3) (a3)--(a4) (a2)--(a5);
\node[vtxs] (b1) at (4.0,-1) {16};
\node[vtxs, fill=black!70, text=white] (b2) at (4.8,-1) {12};
\node[vtxs] (b3) at (4.8,-2) {5};
\draw[aresta] (r4)--(b1) (r4)--(b2) (b2)--(b3);
\node[vtxs] (c1) at (5.7,-1) {20};
\draw[aresta] (r5)--(c1);
\end{tikzpicture}
\end{figura}

\paragraph{Operações.} Descrevemos o heap de máximo; ``aumentar a prioridade'' corresponde à operação \emph{decrease-key} de um heap de mínimo.
\begin{itemize}
  \item \textbf{Inserção}: cria uma árvore de um nó, coloca-a na lista de raízes e atualiza o ponteiro do máximo. Custo real $O(1)$.
  \item \textbf{Fusão}: concatena as duas listas de raízes; o máximo é o maior dos dois. Custo real $O(1)$.
  \item \textbf{Seleção do máximo}: leitura do ponteiro, $O(1)$.
  \item \textbf{Remoção do máximo}: (1) retira a raiz máxima e promove todos os seus filhos a raízes, desmarcados; (2) \textbf{consolida}: percorre as raízes com um vetor $A[0..D]$ indexado pelo grau, onde $D=O(\log n)$ é o grau máximo possível, e, sempre que duas raízes têm o mesmo grau, subordina a de menor prioridade à de maior (o grau aumenta de 1) e repete; (3) reconstrói a lista de raízes a partir de $A$ e encontra o novo máximo. O custo real é $O(D+t)$, onde $t$ é o número de raízes; o amortizado é $O(\log n)$.
  \item \textbf{Aumento de prioridade de $x$}: aumenta a chave; se ela não ficou maior que a do pai, termina. Senão, \textbf{corta} $x$, com sua subárvore, e o coloca na lista de raízes, desmarcado. Em seguida, olhando o pai $y$: se $y$ está \emph{desmarcado}, ele é marcado (a menos que seja raiz) e o processo termina; se $y$ já estava \emph{marcado}, isto é, já tinha perdido um filho, $y$ também é cortado, desmarcado e levado à lista de raízes, e o processo continua com o pai de $y$. São os \textbf{cortes em cascata}. Custo amortizado $O(1)$.
  \item \textbf{Remoção de um nó qualquer}: aumenta a prioridade para $+\infty$ e remove o máximo: $O(\log n)$ amortizado.
  \item \textbf{Diminuição de prioridade}: não é a operação barata desta estrutura; faz-se removendo e reinserindo o nó, em $O(\log n)$ amortizado.
\end{itemize}

\begin{exemplo}[Inserções preguiçosas e uma consolidação]\label{ex:consolida}
Num heap de Fibonacci vazio, inserimos $3,8,5,12,6,9,2$. Cada inserção apenas acrescenta um nó à lista de raízes: são 7 árvores de um nó, e o ponteiro do máximo aponta para 12. Agora removemos o máximo. O 12 sai (não tem filhos) e a consolidação percorre as raízes $3,8,5,6,9,2$, usando o vetor $A$ indexado pelo grau:
\begin{itemize}
  \item $3$: $A[0]=3$. \quad $8$: colide com $A[0]=3$; o link põe 8 por cima, e a árvore $8\to3$ tem grau 1: $A[0]$ fica vazio e $A[1]=8\to3$.
  \item $5$: $A[0]=5$. \quad $6$: colide com 5, formando $6\to5$, de grau 1, que colide com $A[1]=8\to3$; o link põe 8 por cima, formando uma árvore de grau 2 em $A[2]$.
  \item $9$: $A[0]=9$. \quad $2$: colide com 9, formando $9\to2$ em $A[1]$.
\end{itemize}
Restam duas raízes, 8 (grau 2) e 9 (grau 1), e o novo máximo é 9 (Figura~\ref{fig:consolida}). A remoção custou bem mais que uma inserção, mas pagou de uma vez a ``preguiça'' das sete inserções.
\begin{figura}[A consolidação do Exemplo~\ref{ex:consolida}: antes, sete raízes; depois de remover o 12 e consolidar, duas árvores de graus distintos.\label{fig:consolida}]
\begin{tabular}{@{}c@{\hspace{14mm}}c@{}}
\begin{tikzpicture}[x=1.25cm]
\foreach \v [count=\i] in {3,8,5,12,6,9,2} {\node[vtx] (q\i) at (\i,0) {\v};}
\node[vtx, fill=ruim] at (q4) {12};
\foreach \i [evaluate=\i as \j using int(\i+1)] in {1,...,6} {\draw[{Stealth[length=1.5mm]}-{Stealth[length=1.5mm]}, gray] (q\i) -- (q\j);}
\node[font=\scriptsize\sffamily, MidnightBlue] (mx) at (4,0.9) {max};
\draw[seta, MidnightBlue] (mx) -- (q4);
\end{tikzpicture}
&
\begin{forest} bst, baseline [8 [6 [5]] [3]] \end{forest}\quad
\begin{forest} bst, baseline [9,feito [2]] \end{forest}\\[3pt]
\small (a) após 7 inserções & \small (b) após remover o máximo
\end{tabular}
\end{figura}
\end{exemplo}

\paragraph{Regras das marcas.} Todos os nós começam desmarcados. Um nó é marcado quando \emph{perde o primeiro filho} por um corte, e, ao perder o segundo, ele próprio é cortado. \textbf{Raízes nunca estão marcadas} e podem perder quantos filhos for preciso; quando um nó é subordinado a outro, fica desmarcado.

\begin{exemplo}[Cortes em cascata]\label{ex:cascata}
O nó interno $X$ tem pai $Y$, marcado, e avô $Z$, desmarcado, que não é raiz. A prioridade de $X$ aumenta e passa a superar a de $Y$.
\begin{enumerate}
  \item $X$ é cortado e vai para a lista de raízes, desmarcado.
  \item $Y$ acaba de perder o segundo filho: é cortado, desmarcado e levado à lista de raízes.
  \item $Z$, que não é raiz e estava desmarcado, perdeu o seu primeiro filho: passa a ser marcado, e a cascata termina.
\end{enumerate}
Ao final, $X$ e $Y$ estão na lista de raízes, desmarcados, e $Z$ continua no seu lugar, agora marcado: se ele perder mais um filho no futuro, será cortado também. A Figura~\ref{fig:cascata} mostra o exemplo com chaves concretas: a de $X$ sobe de 10 para 60.
\begin{figura}[Cortes em cascata. (a) $Y=30$ está marcado (escuro) e $Z=40$, não; a chave de $X$ passa de 10 a 60. (b) $X$ e $Y$ são cortados, um após o outro, e vão para a lista de raízes desmarcados; $Z$ perde um filho e passa a ser marcado. O novo máximo é 60.\label{fig:cascata}]
\forestset{marc/.style={fill=black!70, text=white}}
\begin{tabular}{@{}c@{\hspace{14mm}}c@{}}
\begin{forest} bst, baseline
[50 [45] [40,label={[font=\scriptsize]right:$Z$} [38] [30,marc,label={[font=\scriptsize]right:$Y$} [25] [60,novo,label={[font=\scriptsize]right:$X$}]]]]
\end{forest}
&
\begin{forest} bst, baseline [50 [45] [40,marc,label={[font=\scriptsize]right:$Z$} [38]]] \end{forest}\quad
\begin{forest} bst, baseline [60,novo] \end{forest}\quad
\begin{forest} bst, baseline [30 [25]] \end{forest}\\[3pt]
\small (a) antes dos cortes & \small (b) depois: três raízes, e $Z$ marcado
\end{tabular}
\end{figura}
\end{exemplo}

\subsection{\texorpdfstring{Por que funciona: o grau máximo é $O(\log n)$}{Por que funciona: o grau máximo é O(log n)}}
\begin{lema}
Seja $x$ um nó e sejam $y_1,\dots,y_k$ os seus filhos, na ordem em que foram ligados a $x$. Então $\operatorname{grau}(y_i)\ge i-2$.
\end{lema}
\begin{proof}
Quando $y_i$ foi ligado a $x$, este já tinha os filhos $y_1,\dots,y_{i-1}$, e portanto grau $\ge i-1$. Só se ligam nós de mesmo grau, logo $y_i$ tinha grau $\ge i-1$. Desde então, $y_i$ perdeu no máximo um filho, pois ao perder o segundo seria cortado. Logo $\operatorname{grau}(y_i)\ge i-2$.
\end{proof}
\begin{teo}
Um nó de grau $k$ tem uma subárvore com pelo menos $F_{k+2}\ge\varphi^k$ nós, onde $F$ é a sequência de Fibonacci e $\varphi=\frac{1+\sqrt5}2$. Logo o grau máximo é $D(n)\le\log_\varphi n=O(\log n)$, e depois de uma consolidação o heap tem $O(\log n)$ árvores.
\end{teo}
\begin{proof}[Esboço]
Seja $s_k$ o menor tamanho possível de uma subárvore cuja raiz tem grau $k$. Pelo lema, $s_k\ge 2+\sum_{i=2}^{k}s_{i-2}$, e por indução $s_k\ge F_{k+2}$, pois $F_{k+2}=1+\sum_{i=0}^{k}F_i$. É daí que vem o nome da estrutura: as menores árvores de posto $0,1,2,3,4,5$ têm $1,2,3,5,8,13$ nós.
\end{proof}

\begin{figura}[As menores árvores possíveis de grau $0$ a $4$ num heap de Fibonacci. Cada uma é uma raiz cujos filhos são as menores árvores de graus $0,0,1,2,\dots,k-2$ (o lema). Os tamanhos $1,2,3,5,8$ são números de Fibonacci.\label{fig:fibmin}]
\forestset{bn/.style={for tree={circle, draw, minimum size=2.8mm, inner sep=0pt, s sep=1.6mm, l sep=1.6mm, edge={thick}}}}
\begin{tabular}{@{}c@{\hspace{9mm}}c@{\hspace{9mm}}c@{\hspace{9mm}}c@{\hspace{9mm}}c@{}}
\begin{forest} bn, baseline [] \end{forest} &
\begin{forest} bn, baseline [ [] ] \end{forest} &
\begin{forest} bn, baseline [ [] [] ] \end{forest} &
\begin{forest} bn, baseline [ [] [] [ [] ] ] \end{forest} &
\begin{forest} bn, baseline [ [] [] [ [] ] [ [] [] ] ] \end{forest}\\[2pt]
\small grau 0: 1 nó & \small grau 1: 2 & \small grau 2: 3 & \small grau 3: 5 & \small grau 4: 8
\end{tabular}
\end{figura}

Sem as marcas e os cortes em cascata, seria possível cortar todos os netos de uma raiz e obter uma árvore de posto $k$ com apenas $k+1$ nós. O posto deixaria de ser $O(\log n)$ e a consolidação ficaria cara. Compare: uma $B_4$ do heap binomial tem 16 nós; a menor árvore de grau 4 de um heap de Fibonacci tem 8; sem marcas, teria só 5 (uma raiz e quatro folhas).

\subsection{Análise amortizada}
Use o potencial $\Phi(H)=t(H)+2\,m(H)$, onde $t$ é o número de raízes e $m$ o número de nós marcados.
\begin{itemize}
  \item \textbf{Inserção}: custo real $O(1)$ e $\Delta\Phi=+1$, logo $O(1)$ amortizado. \textbf{Fusão}: $\Delta\Phi=0$, logo $O(1)$.
  \item \textbf{Remoção do máximo}: custo real $O(D+t)$; depois dela há no máximo $D+1$ raízes, então $\Delta\Phi\le(D+1)-t$, e o custo amortizado é $O(D)=O(\log n)$. O trabalho de consolidar as $t$ raízes acumuladas foi pago pelas inserções e fusões preguiçosas que as criaram.
  \item \textbf{Aumento de prioridade} com $c$ cortes em cascata: custo real $O(c)$; $t$ aumenta de $c$ e $m$ diminui de pelo menos $c-2$, então $\Delta\Phi\le c+2(2-c)=4-c$, e o custo amortizado é $O(1)$. Na linguagem do banqueiro, o corte de um nó marcado é pago pelos créditos deixados pela operação que o marcou.
\end{itemize}

\begin{exemplo}[O potencial nos exemplos anteriores]\label{ex:potencial}
No Exemplo~\ref{ex:consolida}, depois das 7 inserções há $t=7$ raízes e nenhuma marca: $\Phi=7$. A remoção do máximo examinou 6 raízes e fez 4 links, mas terminou com $t=2$, e $\Phi$ caiu de 7 para 2: a queda de 5 unidades ``pagou'' quase todo o trabalho. No Exemplo~\ref{ex:cascata}, houve $c=2$ cortes; $t$ aumentou de 2, e o número de marcas não mudou (a de $Y$ sumiu e a de $Z$ apareceu), logo $\Delta\Phi=2+2\cdot0=2$, e o custo amortizado é $2+2=4$, constante. Cada corte a mais na cascata desmarcaria mais um nó, e o potencial cairia 2 para pagá-lo.
\end{exemplo}

\begin{teo}
A partir de um heap de Fibonacci vazio, qualquer sequência de $a_1$ inserções, $a_2$ remoções do máximo, $a_3$ aumentos de prioridade e $a_4$ fusões custa $O(a_1+a_3+a_4+a_2\log n)$.
\end{teo}

\begin{cuidado}
\begin{itemize}
  \item Os custos $O(1)$ do heap de Fibonacci são \textbf{amortizados}: \emph{uma} remoção do máximo depois de muitas inserções e fusões pode custar $\Theta(n)$, ao consolidar tudo de uma vez. Se cada operação precisa ter tempo previsível, essa estrutura é inadequada.
  \item As constantes são grandes e cada nó ocupa bastante memória (quatro ponteiros, grau e marca). Na prática, a estrutura só compensa com um número muito grande de aumentos de prioridade, como nos algoritmos de Dijkstra e Prim em grafos densos.
  \item Uma árvore de um heap de Fibonacci pode ter altura grande, até $\Theta(n)$ em sequências adversárias. O que é logarítmico é o \textbf{grau}, não a altura.
\end{itemize}
\end{cuidado}

% ------------------------------------------------------------------
\section{Comparação e escolha da estrutura}\label{sec:escolhaheap}
\begin{center}
\small
\begin{tabular}{@{}l>{\centering\arraybackslash}p{3.3cm}>{\centering\arraybackslash}p{4.2cm}>{\centering\arraybackslash}p{3.6cm}@{}}
\toprule
& Heap binário & Heap binomial & Heap de Fibonacci\\
\midrule
Seleção & $O(1)$ & $O(\log n)$ ou $O(1)$ com ponteiro & $O(1)$\\
Inserção & $O(\log n)$ & $O(\log n)$ pior; $O(1)$ amortizado & $O(1)$\am\\
Remoção do máximo & $O(\log n)$ & $O(\log n)$ & $O(\log n)$\am\\
Aumentar prioridade & $O(\log n)$ & $O(\log n)$ & $O(1)$\am\\
Fusão & $O(n)$ & $O(\log n)$ \textbf{pior caso} & $O(1)$\am\\
Construção & $O(n)$ & $O(n)$ & $O(n)$\\
Memória & vetor, sem ponteiros & 3 ponteiros $+$ grau por nó & 4 ponteiros $+$ grau $+$ marca por nó\\
Garantia & pior caso & pior caso & amortizada\\
Simplicidade & muito simples & média & complexa\\
\bottomrule
\end{tabular}
\end{center}

\begin{roteiro}[Como escolher a lista de prioridades]
\begin{enumerate}
  \item Não há fusões? Use o \textbf{heap binário}: é o mais simples, usa menos memória e tem boa localidade de acesso.
  \item Há fusões frequentes e \textbf{cada operação precisa ter tempo previsível} (limite de tempo por operação, sistemas de tempo real)? Use o \textbf{heap binomial}, com fusão $O(\log n)$ no pior caso.
  \item Há uma quantidade muito grande de \textbf{aumentos de prioridade} (ou de fusões) e o que importa é o \textbf{custo total} da execução? Use o \textbf{heap de Fibonacci}.
  \item Há uma restrição forte de \textbf{memória}? Use o heap binário; as outras estruturas pagam vários ponteiros por elemento.
\end{enumerate}
\end{roteiro}

\begin{exemplo}[Aplicando o roteiro]\leavevmode
\begin{itemize}
  \item Uma fila de impressão que só recebe trabalhos e despacha o mais urgente: não há fusões, e o heap binário basta.
  \item Um escalonador de tempo real em que filas de processadores são fundidas quando um processador falha, e cada operação tem prazo: heap binomial.
  \item O algoritmo de Dijkstra num grafo denso, com $n$ remoções e até $m\approx n^2$ aumentos de prioridade: o heap de Fibonacci dá $O(m+n\log n)$, contra $O(m\log n)$ do binário (Capítulo~\ref{cap:dijkstra}).
\end{itemize}
\end{exemplo}

\begin{chave}
Heap binomial $=$ árvores $B_k$ de ordens distintas $=$ bits de $n$; fusão $=$ soma binária, $O(\log n)$ no pior caso; inserção $O(1)$ amortizado. Heap de Fibonacci $=$ preguiça $+$ marcas; tudo $O(1)$ amortizado, exceto remoções, $O(\log n)$ amortizado; grau $\le\log_\varphi n$. Amortizado é garantia para a sequência, não para cada operação, e não é caso médio.
\end{chave}

% ----------------------------------------------------------------
\section{Exercícios}
\excurso
\begin{exercicio}{hb-l7}[Lista]<sec:heapbinomial>
Desenvolva um algoritmo $O(n)$ para construir um heap binomial com $n$ elementos.
\end{exercicio}
\begin{exercicio}{hb-l8}[Lista]<sec:heapbinomial>
Um heap binomial tem $k$ árvores. Que alturas elas devem ter para que a inserção seja o pior caso, isto é, faça o número máximo de passos?
\end{exercicio}
\begin{exercicio}{hb-l10}[Lista]<ex:rembin>
$H$ é um heap binomial com $k$ árvores, e a árvore cuja raiz tem a maior prioridade tem altura $h$. Determine as alturas das árvores de $H$ e o valor de $h$ para que a remoção do máximo faça o número máximo de passos.
\end{exercicio}
\begin{exercicio}{hb-l11}[Lista]<ex:fusaobin>
Prove ou dê contraexemplo: se $H_1$ e $H_2$ têm $k_1$ e $k_2$ árvores e sua fusão $H$ tem $k$ árvores, então $k\ge\min\{k_1,k_2\}$.
\end{exercicio}
\begin{exercicio}{hb-l12}[Lista]<ex:fusaobin>
Prove ou dê contraexemplo: nas mesmas condições, $k\le k_1+k_2$.
\end{exercicio}
\begin{exercicio}{hb-l13}[Lista]<sec:fib>
Determine a altura mínima e a máxima de uma árvore de um heap de Fibonacci em função do posto $k$.
\end{exercicio}
\begin{exercicio}{hb-l14}[Lista]<sec:heapbinomial,sec:fib>
Escreva os algoritmos de seleção, inserção, remoção, aumento e diminuição de prioridade, construção e fusão para heaps binomiais e para heaps de Fibonacci.
\end{exercicio}
\begin{exercicio}{hb-r6}[Monitoria]<sec:binomial>
(a) O que caracteriza uma árvore binomial $B_k$? Quantos nós tem? Qual a relação entre o grau da raiz e $k$? (b) Como um heap binomial representa uma coleção de árvores binomiais? (c) Duas árvores $B_2$ de um heap binomial de mínimo têm raízes de chaves 10 e 15: qual é o resultado do \emph{link}, sua ordem, seu número de nós e o grau da raiz? (d) Como achar o elemento de maior prioridade e qual o custo? (e) Qual o custo da união e por que ela não percorre os $n$ elementos?
\end{exercicio}
\begin{exercicio}{hb-r9}[Monitoria]<sec:fib>
(a) Qual a principal diferença estrutural entre heap de Fibonacci e heap binomial? (b) Quais operações são $O(1)$ amortizado? Quanto custa remover o mínimo (\emph{Extract-Min})? (c) Na \emph{Decrease-Key}$(x,k)$ de um heap de mínimo, por que alterar a chave pode violar a propriedade de heap e o que se faz? (d) Para que servem os cortes em cascata, e por que a operação é $O(1)$ amortizado mesmo fazendo vários cortes?
\end{exercicio}
\begin{exercicio}{hb-m1}[Monitoria]<ex:fusaoformas>
Faça a fusão de $H_1=\{B_0,B_2,B_4\}$ com $H_2=\{B_0,B_1\}$ como uma soma binária: quais subordinações acontecem, em que ordem, e qual o resultado em binário?
\end{exercicio}
\begin{exercicio}{hb-m2}[Monitoria]<ex:cascata>
Num heap de Fibonacci de máximo, o nó interno $X$ tem pai $Y$ marcado e avô $Z$ marcado. A prioridade de $X$ aumenta e viola a ordem. Descreva os cortes em cascata, os nós movidos para a lista de raízes e as mudanças nas marcas.
\end{exercicio}
\begin{exercicio}{hb-m3}[Monitoria]<sec:fib>
Por que não podemos simplesmente cortar um nó do pai sem marcar nem cortar em cascata? O que aconteceria com a relação entre a ordem de uma árvore e o seu número de nós?
\end{exercicio}
\begin{exercicio}{hb-t1}[Avaliação]<sec:escolhaheap>
Um aeroporto mantém uma fila de prioridade por setor de controle de tráfego. Fusões de filas (transferências entre setores, unificação de setores) são muito frequentes, com filas de centenas de milhares de aeronaves; alterações de prioridade são raras; e o tempo de \emph{cada} fusão precisa ser previsível, pois as janelas de processamento são rígidas. Qual heap você escolheria? Justifique e explique por que os outros dois são menos adequados.
\end{exercicio}
\begin{exercicio}{hb-t2}[Avaliação]<sec:escolhaheap>
Uma central de atendimento médico mantém cerca de $2\cdot10^6$ chamadas numa fila de prioridade, com $2\cdot10^6$ inserções, $2\cdot10^6$ remoções da chamada mais urgente e poucas alterações de prioridade. Cada chamada ocupa 16 bytes, endereços ocupam 8 bytes e a estrutura pode usar no máximo $40$ MB $=41\,943\,040$ bytes. Qual estrutura você usaria? Considere a complexidade das operações, a necessidade de fusão, a memória por elemento e a simplicidade.
\end{exercicio}
\begin{exercicio}{hb-t3}[Avaliação]<sec:escolhaheap>
Um supercomputador tem cerca de $5\cdot10^6$ tarefas numa fila de prioridade. Num período típico há $5\cdot10^6$ inserções, $5\cdot10^6$ remoções da tarefa mais prioritária, $8\cdot10^7$ aumentos de prioridade (feitos sem retirar e reinserir a tarefa) e poucas uniões de filas. O custo total da execução é o que importa. Qual heap é mais adequado? Por que os outros dois são menos vantajosos?
\end{exercicio}

\excompl
\begin{exercicio}{hb-vetor}[Complementar]<sec:amort>
Um vetor dinâmico dobra de capacidade quando fica cheio, copiando todos os elementos. Mostre, com o método do potencial e $\Phi=2\cdot(\text{número de elementos})-\text{capacidade}$, que cada inserção custa $O(1)$ amortizado.
\end{exercicio}
\begin{exercicio}{hb-contador}[Complementar]<ex:contador>
Um contador binário de $k$ bits passa a aceitar \emph{incremento} e \emph{decremento}. Mostre uma sequência de $m$ operações que custa $\Theta(mk)$. Por que a análise amortizada do contador só com incrementos não se aplica?
\end{exercicio}
\begin{exercicio}{hb-fila}[Complementar]<sec:amort>
Implemente uma fila com duas pilhas (uma de entrada e outra de saída) e mostre que cada operação custa $O(1)$ amortizado.
\end{exercicio}
\begin{exercicio}{hb-forma}[Complementar]<sec:heapbinomial>
Quais árvores binomiais compõem um heap binomial com 100 elementos? E depois de mais uma inserção? E depois de mais quatro inserções a partir de 100? Quantos \emph{links} cada uma dessas inserções faz?
\end{exercicio}
\begin{exercicio}{hb-fibbin}[Complementar]<sec:fib>
Mostre que, se um heap de Fibonacci sofre apenas inserções, fusões e remoções do máximo (sem alterações de prioridade), toda árvore dele é uma árvore binomial (não ordenada) logo após cada consolidação.
\end{exercicio}
\begin{exercicio}{hb-fibcaro}[Complementar]<ex:consolida>
Dê uma sequência de operações, a partir de um heap de Fibonacci vazio, em que \emph{uma única} remoção do máximo tem custo real $\Theta(n)$. Por que isso não contradiz o custo amortizado $O(\log n)$?
\end{exercicio}
\begin{exercicio}{hb-binsim}[Complementar]<ex:fusaobin>
Insira $4,9,2,7,5$, nessa ordem, num heap binomial de máximo inicialmente vazio. Mostre as árvores depois de cada inserção e conte os \emph{links}. Compare o número de links de cada inserção com o número de ``vai um'' ao somar 1 em binário.
\end{exercicio}
\begin{exercicio}{hb-fibsim}[Complementar]<ex:consolida>
Num heap de Fibonacci de máximo vazio, insira $6,2,9,4,8,1,7$ e depois remova o máximo, consolidando as raízes na ordem em que foram inseridas, como no Exemplo~\ref{ex:consolida}. Quais árvores restam? Quantos links são feitos? Quanto vale o potencial $\Phi=t+2m$ antes e depois da remoção?
\end{exercicio}
